\documentclass[10pt,a4paper]{article}
\usepackage[utf8]{inputenc}
\usepackage[T1]{fontenc}
\usepackage[english]{babel}
\usepackage{geometry}
\geometry{margin=1.8cm}
\usepackage{enumitem}
\usepackage{booktabs}
\usepackage{fancyhdr}
\usepackage{lastpage}
\usepackage{tcolorbox}

\pagestyle{fancy}
\fancyhf{}
\rfoot{\thepage/\pageref{LastPage}}
\lhead{\small TOEFL ITP 2026 --- Day 7: Review}
\rhead{\small MPrometeo}
\renewcommand{\headrulewidth}{0.4pt}
\setlength{\parindent}{0pt}
\setlength{\parskip}{5pt}

\begin{document}

\begin{center}
{\LARGE \textbf{DAY 7 --- Thursday Mar 5}}\\[4pt]
{\large Review \& Rest}\\[2pt]
{\normalsize Theme: Maple vs Mathematica + Grammar Cheat Sheet}\\[2pt]
\rule{\textwidth}{0.8pt}
\end{center}

% ============================================================
\section*{Reading: Maple vs Mathematica ($\sim$230 words)}
% ============================================================

\begin{tcolorbox}[colback=white, colframe=black]
\small
The debate between Maple and Mathematica has persisted for over three
decades in academic mathematics and physics departments. Both are
computer algebra systems (CAS) capable of symbolic computation,
numerical analysis, and visualization, yet they differ in philosophy
and user experience.

Mathematica, developed by Wolfram Research, emphasizes a unified
computational language. Its creator, Stephen Wolfram, envisions it
as a universal knowledge engine. The software excels at integration
between symbolic math, data analysis, machine learning, and
publication-quality graphics. Its notebook interface influenced the
design of Jupyter notebooks. However, critics describe its syntax as
opaque and its licensing as expensive --- a perpetual license can
cost over \$3,000 for an individual.

Maple, developed by Maplesoft (originally from the University of
Waterloo), takes a more traditional mathematical approach. Many
mathematicians find its syntax closer to how they write equations on
paper. Maple's strength lies in pure symbolic computation ---
particularly differential equations and algebraic manipulation.
Its pricing, while still substantial, tends to be lower than
Mathematica's for academic users.

In practice, the choice often depends on department tradition rather
than objective superiority. Physics departments tend toward
Mathematica; engineering departments lean toward Maple or MATLAB.
Meanwhile, the rise of free alternatives --- SageMath, SymPy, and
Julia --- is gradually eroding the dominance of both. For a graduate
student on a budget, the ``best'' CAS may simply be the one you can
afford.
\end{tcolorbox}

\subsection*{Questions}

\begin{enumerate}[label=\textbf{\arabic*.}]
\item The passage is mainly about:
\begin{enumerate}[label=(\Alph*)]
\item Why Mathematica is superior
\item A comparison of two major computer algebra systems
\item The history of programming
\item Why free software is always better
\end{enumerate}

\item ``Opaque'' in describing Mathematica's syntax means:
\begin{enumerate}[label=(\Alph*)]
\item transparent and clear \item difficult to understand
\item colorful \item fast
\end{enumerate}

\item According to the passage, department choice depends on:
\begin{enumerate}[label=(\Alph*)]
\item Objective performance benchmarks
\item Tradition and discipline
\item Student votes
\item Government regulations
\end{enumerate}

\item ``Eroding the dominance'' means:
\begin{enumerate}[label=(\Alph*)]
\item Supporting the leaders
\item Gradually weakening the market control of both
\item Destroying all software
\item Making them cheaper
\end{enumerate}

\item The author's tone about free alternatives is:
\begin{enumerate}[label=(\Alph*)]
\item Dismissive \item Cautiously optimistic
\item Angry \item Indifferent
\end{enumerate}
\end{enumerate}

\textbf{Answers:} 1-B, 2-B, 3-B, 4-B, 5-B

% ============================================================
\section*{Grammar Mega-Review: 15 TOEFL Patterns Cheat Sheet}
% ============================================================

\textbf{Memorize these. They cover 80\% of Structure questions.}

\begin{center}
\small
\begin{tabular}{clll}
\toprule
\# & \textbf{Pattern} & \textbf{Error} & \textbf{Correct} \\
\midrule
1 & SVA & The list of items are & \textbf{is} \\
2 & Parallel & likes reading, to swim & reading, \textbf{swimming} \\
3 & Word Form & grew rapid & grew \textbf{rapidly} \\
4 & Passive & The bridge built in 1900 & \textbf{was built} \\
5 & Relative & The man which came & \textbf{who} came \\
6 & Gerund/Inf & avoids to eat & avoids \textbf{eating} \\
7 & Comparative & more easier & \textbf{easier} \\
8 & Articles & She is professor & She is \textbf{a} professor \\
9 & Conditional & If I would know & If I \textbf{knew} \\
10 & Inversion & Not only she is & Not only \textbf{is she} \\
11 & Reduced & man stood there & man \textbf{standing} there \\
12 & Subjunctive & suggested she goes & suggested she \textbf{go} \\
13 & Despite/Although & Despite it rains & \textbf{Although} it rains \\
14 & Number/Amount & amount of students & \textbf{number} of students \\
15 & Its/It's & it's interface & \textbf{its} interface \\
\bottomrule
\end{tabular}
\end{center}

% ============================================================
\section*{Week 1 Reflection Template}
% ============================================================

\begin{tcolorbox}[colback=white, colframe=black]
Fill in after completing Days 1--7:

\begin{tabular}{ll}
\textbf{My 3 weakest areas:} & 1. \rule{4cm}{0.4pt} \\
& 2. \rule{4cm}{0.4pt} \\
& 3. \rule{4cm}{0.4pt} \\[6pt]
\textbf{Most common error type:} & \rule{5cm}{0.4pt} \\[6pt]
\textbf{Estimated score now:} & L: \rule{1cm}{0.4pt}
S: \rule{1cm}{0.4pt} R: \rule{1cm}{0.4pt} \\[6pt]
\textbf{What to change next week:} & \rule{5cm}{0.4pt} \\
\end{tabular}
\end{tcolorbox}

\section*{Shadowing (Gemini Live)}

\begin{quote}\small
``The debate between Maple and Mathematica has persisted for over
three decades. Both are computer algebra systems capable of symbolic
computation, numerical analysis, and visualization, yet they differ
in philosophy and user experience. For a graduate student on a
budget, the best system may simply be the one you can afford.''
\end{quote}

\textbf{Targets:} \textbf{persisted} (per-SIS-ted), \textbf{algebra}
(AL-je-bra), \textbf{visualization} (vizh-u-al-i-ZAY-shun),
\textbf{philosophy} (fi-LOS-o-fee).

\end{document}